\documentclass[11pt,a4paper]{article}

% Pakete
\usepackage[utf8]{inputenc}
\usepackage[T1]{fontenc}
\usepackage[ngerman]{babel}
\usepackage{geometry}
\usepackage{graphicx}
\usepackage{xcolor}
\usepackage{booktabs}
\usepackage{longtable}
\usepackage{enumitem}
\usepackage{fancyhdr}
\usepackage{titlesec}
\usepackage{hyperref}
\usepackage{tcolorbox}
\usepackage{array}
\usepackage{colortbl}
\usepackage{tikz}
\usepackage{amsmath}

% Seitenränder
\geometry{left=2.5cm, right=2.5cm, top=3cm, bottom=3cm}

% FehrAdvice Farben
\definecolor{fadarkblue}{RGB}{2,64,121}
\definecolor{falightblue}{RGB}{84,158,222}
\definecolor{fadarkgray}{RGB}{37,33,42}
\definecolor{falightgray}{RGB}{243,245,247}
\definecolor{faorange}{RGB}{222,157,62}
\definecolor{fared}{RGB}{200,60,60}
\definecolor{fagreen}{RGB}{60,160,60}
\definecolor{fapurple}{RGB}{120,80,160}

% Hyperlinks
\hypersetup{
    colorlinks=true,
    linkcolor=fadarkblue,
    urlcolor=falightblue
}

% Kopf- und Fusszeile
\pagestyle{fancy}
\fancyhf{}
\fancyhead[L]{\textcolor{fadarkgray}{\small Behavioral Impact Cockpit -- Phase 2}}
\fancyhead[R]{\textcolor{fadarkgray}{\small BFE-2026-002-COCKPIT-P2}}
\fancyfoot[C]{\textcolor{fadarkgray}{\thepage}}
\fancyfoot[R]{\textcolor{fadarkgray}{\small Version 1.0}}
\renewcommand{\headrulewidth}{0.5pt}
\renewcommand{\footrulewidth}{0.5pt}

% Überschriften-Formatierung
\titleformat{\section}{\Large\bfseries\color{fadarkblue}}{\thesection}{1em}{}
\titleformat{\subsection}{\large\bfseries\color{fadarkblue}}{\thesubsection}{1em}{}
\titleformat{\subsubsection}{\normalsize\bfseries\color{fadarkgray}}{\thesubsubsection}{1em}{}

% Box-Stile
\tcbset{
    referenzbox/.style={
        colback=falightblue!10,
        colframe=falightblue,
        fonttitle=\bfseries,
        title=#1
    },
    zielbox/.style={
        colback=fagreen!10,
        colframe=fagreen,
        fonttitle=\bfseries,
        title=#1
    },
    risikobox/.style={
        colback=faorange!10,
        colframe=faorange,
        fonttitle=\bfseries,
        title=#1
    },
    infobox/.style={
        colback=falightgray,
        colframe=fadarkgray,
        fonttitle=\bfseries,
        title=#1
    }
}

\begin{document}

% Titelseite
\begin{titlepage}
    \centering
    \vspace*{2cm}

    {\Huge\bfseries\textcolor{fadarkblue}{Behavioral Impact Cockpit}}\\[0.5cm]
    {\LARGE\textcolor{fadarkblue}{Phase 2: Skalierung \& Optimierung}}\\[2cm]

    {\Large Projektplan}\\[0.5cm]
    {\large Version 1.0}\\[3cm]

    \begin{tcolorbox}[referenzbox={Bezug zu Phase 1}]
        \textbf{Vorgängerprojekt:} BFE-2026-001-COCKPIT\\
        \textbf{Dokumentation:} projektplan\_vollstaendig.yaml\\
        \textbf{Laufzeit Phase 1:} 19.12.2025 -- 01.03.2026\\
        \textbf{Status:} Abgeschlossen
    \end{tcolorbox}

    \vfill

    \begin{tabular}{ll}
        \textbf{Projekt-ID:} & BFE-2026-002-COCKPIT-P2\\
        \textbf{Auftraggeber:} & Bundesamt für Energie (BFE)\\
        \textbf{Programm:} & Energie Schweiz für Private (ECHfP)\\
        \textbf{Laufzeit Phase 2:} & 01.03.2026 -- 30.09.2026 (7 Monate)\\
        \textbf{Budget:} & CHF 145'000\\
        \textbf{Erstellt:} & 02.02.2026\\
    \end{tabular}

    \vspace{1cm}
    {\small\textcolor{fadarkgray}{FehrAdvice \& Partners AG}}
\end{titlepage}

% Inhaltsverzeichnis
\tableofcontents
\newpage

% =============================================================================
\section{Bezug zum bisherigen Projektplan (Phase 1)}
% =============================================================================

\begin{tcolorbox}[referenzbox={Phase 1 -- Zusammenfassung}]
Phase 2 baut direkt auf den Ergebnissen von Phase 1 auf. Alle Referenzen beziehen sich auf das Dokument \texttt{projektplan\_vollstaendig.yaml}.
\end{tcolorbox}

\subsection{Abgeschlossene Meilensteine Phase 1}

\begin{tabular}{>{\columncolor{falightgray}}l l l}
\toprule
\textbf{ID} & \textbf{Meilenstein} & \textbf{Datum}\\
\midrule
MS1 & KPI-Architektur finalisiert & 19.12.2025\\
MS2 & Fragebogen abgenommen & 15.01.2026\\
MS3 & Dashboard-Prototyp fertig & 31.01.2026\\
MS4 & Erstmessung durchgeführt & 07.02.2026\\
MS5 & Erste Auswertung geliefert & 15.02.2026\\
MS6 & Cockpit übergeben & 01.03.2026\\
\bottomrule
\end{tabular}

\subsection{Ergebnisse Phase 1}

Die folgenden Ergebnisse bilden die Grundlage für Phase 2:

\begin{itemize}[leftmargin=*]
    \item Validiertes KPI-Framework (A-W-I-T + Gesamt-Score)
    \item Getesteter Fragebogen (n = 1'500 Befragte)
    \item Funktionales Dashboard mit Echtzeit-Daten
    \item Baseline-Messung aller KPIs
    \item Dokumentierte Berechnungslogik
    \item Geschultes BFE-Team
\end{itemize}

\subsection{Baseline-Werte und Ziele Phase 2}

\begin{tabular}{l c c}
\toprule
\textbf{KPI} & \textbf{Baseline (Phase 1)} & \textbf{Ziel Phase 2}\\
\midrule
\rowcolor{falightgray} Awareness (A) & [Nach Erstmessung] & +10 Prozentpunkte\\
Willingness (W) & [Nach Erstmessung] & +15 Prozentpunkte\\
\rowcolor{falightgray} Impact (I) & [Nach Erstmessung] & +8 Prozentpunkte\\
Trust (T) & [Nach Erstmessung] & +5 Prozentpunkte\\
\rowcolor{falightgray} \textbf{Gesamt-Score} & [Nach Erstmessung] & \textbf{+10 Prozentpunkte}\\
\bottomrule
\end{tabular}

\newpage

% =============================================================================
\section{Phase 2 -- Übersicht}
% =============================================================================

\subsection{Zielsetzung}

\begin{tcolorbox}[zielbox={Hauptziel Phase 2}]
Transformation des Behavioral Impact Cockpits von einem Pilot-Dashboard zu einem \textbf{skalierbaren, automatisierten Wirkungsmessungssystem} für alle BFE-Kampagnen im Bereich Energie Schweiz.
\end{tcolorbox}

\subsection{Teilziele}

\begin{tabular}{c l p{7cm} p{4cm}}
\toprule
\textbf{ID} & \textbf{Ziel} & \textbf{Beschreibung} & \textbf{Messbar}\\
\midrule
\rowcolor{falightgray} Z1 & Longitudinale Analyse & 3 Tracking-Wellen zur KPI-Entwicklung & 3 Zyklen mit n$\geq$1'500\\
Z2 & Kampagnen-Erweiterung & Integration von 2 weiteren BFE-Kampagnen & 2 neue Module live\\
\rowcolor{falightgray} Z3 & Automatisierung & Automatische Datenerfassung \& Reports & Manuelle Aufwände --70\%\\
Z4 & Interventions-Integration & Verknüpfung Interventionen $\leftrightarrow$ KPIs & 3 Interventionen messbar\\
\rowcolor{falightgray} Z5 & Wissenstransfer & Befähigung des BFE-Teams & Eigenständige Nutzung\\
\bottomrule
\end{tabular}

\subsection{Erfolgskriterien (SMART)}

\subsubsection{Primäre Erfolgskriterien}

\begin{tabular}{c l l l}
\toprule
\textbf{ID} & \textbf{Kriterium} & \textbf{Zielwert} & \textbf{Messmethode}\\
\midrule
\rowcolor{falightgray} E1 & KPI-Verbesserung & Gesamt-Score +10 PP & Welle 1 vs. Welle 4\\
E2 & Skalierung erreicht & 3 Kampagnen integriert & Anzahl Module\\
\rowcolor{falightgray} E3 & Automatisierungsgrad & $\geq$70\% automatisiert & Zeit pro Report\\
\bottomrule
\end{tabular}

\subsubsection{Sekundäre Erfolgskriterien}

\begin{tabular}{c l l l}
\toprule
\textbf{ID} & \textbf{Kriterium} & \textbf{Zielwert} & \textbf{Messmethode}\\
\midrule
\rowcolor{falightgray} E4 & Nutzerzufriedenheit & $\geq$4.0 / 5.0 & Feedback-Befragung\\
E5 & Budget-Einhaltung & $\leq$105\% Budget & Kostencontrolling\\
\bottomrule
\end{tabular}

\newpage

% =============================================================================
\section{Budget \& Ressourcen}
% =============================================================================

\subsection{Budgetübersicht}

\begin{center}
\begin{tikzpicture}
    % Balkendiagramm
    \fill[fadarkblue] (0,0) rectangle (6.5,0.8);
    \fill[falightblue] (0,1.2) rectangle (4.5,2.0);
    \fill[faorange] (0,2.4) rectangle (1.5,3.2);
    \fill[fagreen] (0,3.6) rectangle (1.8,4.4);
    \fill[fadarkgray] (0,4.8) rectangle (0.2,5.6);

    % Labels
    \node[right] at (6.7,0.4) {Personal intern: CHF 65'000 (45\%)};
    \node[right] at (4.7,1.6) {Personal extern: CHF 45'000 (31\%)};
    \node[right] at (1.7,2.8) {Infrastruktur: CHF 15'000 (10\%)};
    \node[right] at (2.0,4.0) {Erhebung: CHF 18'000 (12\%)};
    \node[right] at (0.4,5.2) {Reserve: CHF 2'000 (2\%)};
\end{tikzpicture}
\end{center}

\begin{tcolorbox}[infobox={Gesamtbudget Phase 2}]
\centering
\textbf{\Large CHF 145'000}\\[0.5em]
\small (Phase 1: CHF 85'000 $\rightarrow$ +71\% für erweiterten Scope)
\end{tcolorbox}

\subsection{Budgetaufschlüsselung}

\begin{tabular}{l r p{6cm}}
\toprule
\textbf{Position} & \textbf{Betrag} & \textbf{Details}\\
\midrule
\rowcolor{falightgray} Personal intern & CHF 65'000 & Projektleitung, Analysten (ca. 800h)\\
Personal extern & CHF 45'000 & FehrAdvice Beratung, EBF-Integration\\
\rowcolor{falightgray} Infrastruktur & CHF 15'000 & Dashboard-Hosting, API-Erweiterungen\\
Erhebung & CHF 18'000 & 3 Tracking-Wellen à 1'500 Befragte\\
\rowcolor{falightgray} Reserve & CHF 2'000 & Unvorhergesehenes (ca. 1.4\%)\\
\midrule
\textbf{Total} & \textbf{CHF 145'000} & \\
\bottomrule
\end{tabular}

\newpage

% =============================================================================
\section{Arbeitspakete Phase 2}
% =============================================================================

\subsection{Übersicht}

\begin{center}
\begin{tikzpicture}[scale=0.9]
    % Timeline
    \draw[thick, fadarkblue] (0,0) -- (14,0);

    % Monate
    \foreach \x/\m in {0/Mrz, 2/Apr, 4/Mai, 6/Jun, 8/Jul, 10/Aug, 12/Sep} {
        \draw[fadarkblue] (\x,-0.2) -- (\x,0.2);
        \node[below] at (\x,-0.3) {\small \m};
    }

    % Arbeitspakete
    \fill[fadarkblue!60] (0,1) rectangle (1.5,1.5) node[midway, white, font=\tiny] {AP12};
    \fill[falightblue!80] (2,1) rectangle (3.5,1.5) node[midway, white, font=\tiny] {AP13};
    \fill[fagreen!70] (1.5,2) rectangle (5.5,2.5) node[midway, white, font=\tiny] {AP14};
    \fill[fagreen!70] (4,3) rectangle (7,3.5) node[midway, white, font=\tiny] {AP15};
    \fill[faorange!70] (3,4) rectangle (6.5,4.5) node[midway, white, font=\tiny] {AP16};
    \fill[falightblue!80] (6.5,1) rectangle (8,1.5) node[midway, white, font=\tiny] {AP17};
    \fill[fapurple!70] (5.5,2) rectangle (9,2.5) node[midway, white, font=\tiny] {AP18};
    \fill[falightblue!80] (10,1) rectangle (11.5,1.5) node[midway, white, font=\tiny] {AP19};
    \fill[fadarkblue!60] (11.5,1) rectangle (14,1.5) node[midway, white, font=\tiny] {AP20};

    % Legende
    \node[right, font=\small] at (14.5,1.25) {Wellen/Abschluss};
    \node[right, font=\small] at (14.5,2.25) {Kampagnen/Interv.};
    \node[right, font=\small] at (14.5,3.25) {Module};
    \node[right, font=\small] at (14.5,4.25) {Automatisierung};
\end{tikzpicture}
\end{center}

\subsection{AP 12 -- Baseline-Review \& Zieldefinition}

\begin{tabular}{ll}
\textbf{Verantwortlich:} & Lucas, Manuel\\
\textbf{Zeitraum:} & 01.03.2026 -- 15.03.2026\\
\textbf{Bezug Phase 1:} & MS6, AP11\\
\end{tabular}

\textbf{Aufgaben:}
\begin{enumerate}[label=AP12.\arabic*]
    \item Analyse der Erstmessung aus Phase 1
    \item Definition der KPI-Zielwerte für Phase 2
    \item Identifikation von Verbesserungspotentialen im Cockpit
    \item Priorisierung der Phase-2-Massnahmen
\end{enumerate}

\textbf{Ergebnis:} Validierte Baseline, definierte Zielwerte, priorisierte Roadmap

\subsection{AP 13 -- Tracking-Welle 2}

\begin{tabular}{ll}
\textbf{Verantwortlich:} & Andrea, Panel-Partner\\
\textbf{Zeitraum:} & 01.04.2026 -- 15.04.2026\\
\textbf{Bezug Phase 1:} & AP4, AP5\\
\end{tabular}

\textbf{Aufgaben:}
\begin{enumerate}[label=AP13.\arabic*]
    \item Fragebogen-Update basierend auf Phase-1-Learnings
    \item Feldphase Welle 2 (n = 1'500)
    \item Datenaufbereitung und Qualitätsprüfung
    \item KPI-Berechnung und Dashboard-Update
\end{enumerate}

\textbf{Ergebnis:} Welle-2-Daten im Cockpit, erste Trendanalyse

\subsection{AP 14 -- Kampagnen-Modul 2: Heizungsersatz}

\begin{tabular}{ll}
\textbf{Verantwortlich:} & Manuel, BFE-Fachperson\\
\textbf{Zeitraum:} & 15.03.2026 -- 15.05.2026\\
\textbf{Bezug Phase 1:} & AP1, AP2, AP3\\
\end{tabular}

\textbf{Aufgaben:}
\begin{enumerate}[label=AP14.\arabic*]
    \item KPI-Architektur für Heizungsersatz-Kampagne
    \item Fragebogen-Modul entwickeln
    \item Dashboard-Erweiterung
    \item Integration und Test
\end{enumerate}

\textbf{Ergebnis:} Heizungsersatz-Modul live im Cockpit

\subsection{AP 15 -- Kampagnen-Modul 3: Solarenergie}

\begin{tabular}{ll}
\textbf{Verantwortlich:} & Manuel, BFE-Fachperson\\
\textbf{Zeitraum:} & 01.05.2026 -- 30.06.2026\\
\textbf{Bezug Phase 1:} & AP1, AP2, AP3\\
\end{tabular}

\textbf{Aufgaben:}
\begin{enumerate}[label=AP15.\arabic*]
    \item KPI-Architektur für Solarenergie-Kampagne
    \item Fragebogen-Modul entwickeln
    \item Dashboard-Erweiterung
    \item Integration und Test
\end{enumerate}

\textbf{Ergebnis:} Solarenergie-Modul live im Cockpit

\subsection{AP 16 -- Automatisierung \& API-Integration}

\begin{tabular}{ll}
\textbf{Verantwortlich:} & Tech-Lead, Manuel\\
\textbf{Zeitraum:} & 15.04.2026 -- 15.06.2026\\
\textbf{Bezug Phase 1:} & AP3, AP9\\
\end{tabular}

\textbf{Aufgaben:}
\begin{enumerate}[label=AP16.\arabic*]
    \item Automatische Datenimport-Pipeline entwickeln
    \item Scheduled Report-Generierung implementieren
    \item Alert-System für KPI-Abweichungen
    \item API-Dokumentation erstellen
\end{enumerate}

\textbf{Ergebnis:} Automatisierte Pipeline, 70\% weniger manuelle Arbeit

\subsection{AP 17 -- Tracking-Welle 3}

\begin{tabular}{ll}
\textbf{Verantwortlich:} & Andrea, Panel-Partner\\
\textbf{Zeitraum:} & 15.06.2026 -- 30.06.2026\\
\textbf{Bezug Phase 1:} & AP4, AP5\\
\end{tabular}

\textbf{Ergebnis:} Welle-3-Daten, Trendreport über 3 Wellen

\subsection{AP 18 -- Interventions-Integration}

\begin{tabular}{ll}
\textbf{Verantwortlich:} & Lucas, FehrAdvice\\
\textbf{Zeitraum:} & 15.05.2026 -- 31.07.2026\\
\textbf{Bezug Phase 1:} & AP6 (BCJ)\\
\end{tabular}

\textbf{Aufgaben:}
\begin{enumerate}[label=AP18.\arabic*]
    \item Interventions-Katalog mit 10C-Dimensionen erstellen
    \item Verknüpfung Interventionen $\leftrightarrow$ KPIs im Dashboard
    \item A/B-Test-Tracking implementieren
    \item Wirkungsnachweis für 3 Interventionen
\end{enumerate}

\textbf{Ergebnis:} Interventions-Modul im Cockpit, messbare Wirkungsnachweise

\subsection{AP 19 -- Tracking-Welle 4 (Abschluss)}

\begin{tabular}{ll}
\textbf{Verantwortlich:} & Andrea, Panel-Partner\\
\textbf{Zeitraum:} & 15.08.2026 -- 31.08.2026\\
\end{tabular}

\textbf{Ergebnis:} Finale KPI-Entwicklung, Wirkungsnachweis Phase 2

\subsection{AP 20 -- Abschluss \& Übergabe}

\begin{tabular}{ll}
\textbf{Verantwortlich:} & Lucas, Manuel, BFE-Projektleitung\\
\textbf{Zeitraum:} & 01.09.2026 -- 30.09.2026\\
\textbf{Bezug Phase 1:} & AP10, AP11\\
\end{tabular}

\textbf{Aufgaben:}
\begin{enumerate}[label=AP20.\arabic*]
    \item Abschlussbericht Phase 2 erstellen
    \item Dokumentation aller Erweiterungen
    \item Schulung BFE-Team (erweiterte Funktionen)
    \item Empfehlungen für Phase 3 / Regelbetrieb
    \item Projekt-Abnahme durch BFE
\end{enumerate}

\textbf{Ergebnis:} Abgenommenes, skaliertes Cockpit im Regelbetrieb

\newpage

% =============================================================================
\section{Meilensteine Phase 2}
% =============================================================================

\begin{longtable}{c l l p{5cm} l}
\toprule
\textbf{ID} & \textbf{Meilenstein} & \textbf{Datum} & \textbf{Deliverables} & \textbf{Abnahme}\\
\midrule
\endhead

\rowcolor{falightgray}
MS7 & Phase-2-Kickoff & 15.03.2026 &
    \begin{minipage}[t]{5cm}
    \vspace{0.1cm}
    \begin{itemize}[leftmargin=*, nosep]
        \item Validierte Baseline
        \item Priorisierte Roadmap
    \end{itemize}
    \vspace{0.1cm}
    \end{minipage} & BFE-PL\\

MS8 & Welle 2 abgeschlossen & 15.04.2026 &
    \begin{minipage}[t]{5cm}
    \vspace{0.1cm}
    \begin{itemize}[leftmargin=*, nosep]
        \item Welle-2-Daten
        \item Erste Trendanalyse
    \end{itemize}
    \vspace{0.1cm}
    \end{minipage} & BFE-PL\\

\rowcolor{falightgray}
MS9 & Kampagnen-Erweiterung & 30.06.2026 &
    \begin{minipage}[t]{5cm}
    \vspace{0.1cm}
    \begin{itemize}[leftmargin=*, nosep]
        \item Heizungsersatz-Modul
        \item Solarenergie-Modul
    \end{itemize}
    \vspace{0.1cm}
    \end{minipage} & BFE-FB\\

MS10 & Automatisierung & 15.06.2026 &
    \begin{minipage}[t]{5cm}
    \vspace{0.1cm}
    \begin{itemize}[leftmargin=*, nosep]
        \item Auto-Pipeline
        \item Scheduled Reports
        \item Alert-System
    \end{itemize}
    \vspace{0.1cm}
    \end{minipage} & Tech-Rev\\

\rowcolor{falightgray}
MS11 & Interventions-Nachweis & 31.07.2026 &
    \begin{minipage}[t]{5cm}
    \vspace{0.1cm}
    \begin{itemize}[leftmargin=*, nosep]
        \item Welle-3-Trend
        \item 3 Interventionen messbar
    \end{itemize}
    \vspace{0.1cm}
    \end{minipage} & BFE-PL, FA\\

MS12 & \textbf{Phase-2-Abschluss} & 30.09.2026 &
    \begin{minipage}[t]{5cm}
    \vspace{0.1cm}
    \begin{itemize}[leftmargin=*, nosep]
        \item Finale Messung
        \item Abschlussbericht
        \item Regelbetrieb
    \end{itemize}
    \vspace{0.1cm}
    \end{minipage} & \textbf{BFE-GL}\\

\bottomrule
\end{longtable}

\newpage

% =============================================================================
\section{Risikomanagement}
% =============================================================================

\begin{tcolorbox}[risikobox={Risiko-Register Phase 2}]
5 identifizierte Risiken mit Massnahmen (fortlaufende Nummerierung ab Phase 1)
\end{tcolorbox}

\begin{longtable}{c p{4cm} c c p{4cm}}
\toprule
\textbf{ID} & \textbf{Risiko} & \textbf{W} & \textbf{A} & \textbf{Massnahmen}\\
\midrule
\endhead

\rowcolor{faorange!20}
R6 & Panel-Kapazitäten bei 4 Wellen & M & M &
    \begin{minipage}[t]{4cm}
    \vspace{0.1cm}
    \begin{itemize}[leftmargin=*, nosep, font=\small]
        \item Frühzeitige Buchung
        \item Backup-Panel
    \end{itemize}
    \vspace{0.1cm}
    \end{minipage}\\

\rowcolor{fared!20}
R7 & Technische Komplexität Automatisierung & M & H &
    \begin{minipage}[t]{4cm}
    \vspace{0.1cm}
    \begin{itemize}[leftmargin=*, nosep, font=\small]
        \item Agile Sprints
        \item Frühes MVP
    \end{itemize}
    \vspace{0.1cm}
    \end{minipage}\\

\rowcolor{fared!20}
R8 & Unzureichende KPI-Verbesserung & N & H &
    \begin{minipage}[t]{4cm}
    \vspace{0.1cm}
    \begin{itemize}[leftmargin=*, nosep, font=\small]
        \item EBF-Interventionen
        \item Kontinuierliches Monitoring
    \end{itemize}
    \vspace{0.1cm}
    \end{minipage}\\

\rowcolor{faorange!20}
R9 & Ressourcenengpässe BFE-Fachbereiche & M & M &
    \begin{minipage}[t]{4cm}
    \vspace{0.1cm}
    \begin{itemize}[leftmargin=*, nosep, font=\small]
        \item Feste Ansprechpartner
        \item Klare Zeitfenster
    \end{itemize}
    \vspace{0.1cm}
    \end{minipage}\\

\rowcolor{faorange!20}
R10 & Budget-Überschreitung (Scope Creep) & M & M &
    \begin{minipage}[t]{4cm}
    \vspace{0.1cm}
    \begin{itemize}[leftmargin=*, nosep, font=\small]
        \item Strikte Scope-Definition
        \item Change-Request-Prozess
    \end{itemize}
    \vspace{0.1cm}
    \end{minipage}\\

\bottomrule
\end{longtable}

{\small W = Wahrscheinlichkeit (N/M/H), A = Auswirkung (N/M/H)}

\newpage

% =============================================================================
\section{Governance \& Kommunikation}
% =============================================================================

\subsection{Gremienstruktur}

\begin{tabular}{l l l p{4cm}}
\toprule
\textbf{Gremium} & \textbf{Mitglieder} & \textbf{Rhythmus} & \textbf{Entscheidungen}\\
\midrule
\rowcolor{falightgray}
Steuerungsausschuss & BFE-GL, FA-Partner, PL & Alle 2 Monate & Budget, Scope, Eskalationen\\
Projektteam & Lucas, Manuel, Andrea, BFE & Wöchentlich & Operative Umsetzung\\
\bottomrule
\end{tabular}

\subsection{Eskalationsstufen}

\begin{center}
\begin{tikzpicture}
    % Stufe 1
    \fill[fagreen!60] (0,0) rectangle (4,1);
    \node[white, font=\small\bfseries] at (2,0.5) {Stufe 1: Projektleitung};

    % Stufe 2
    \fill[faorange!70] (0,1.2) rectangle (4,2.2);
    \node[white, font=\small\bfseries] at (2,1.7) {Stufe 2: Steuerungsausschuss};

    % Stufe 3
    \fill[fared!70] (0,2.4) rectangle (4,3.4);
    \node[white, font=\small\bfseries] at (2,2.9) {Stufe 3: BFE-Geschäftsleitung};

    % Beschreibungen
    \node[right, font=\small] at (4.3,0.5) {Verzögerungen $<$2 Wochen, Budget $<$10\%};
    \node[right, font=\small] at (4.3,1.7) {Verzögerungen $>$2 Wochen, Budget $>$10\%};
    \node[right, font=\small] at (4.3,2.9) {Projektabbruch, Richtungsänderung};
\end{tikzpicture}
\end{center}

\subsection{Regelkommunikation}

\begin{tabular}{l l l l}
\toprule
\textbf{Typ} & \textbf{Rhythmus} & \textbf{Empfänger} & \textbf{Format}\\
\midrule
\rowcolor{falightgray}
Statusbericht & Monatlich & Steuerungsausschuss & PPT, max. 10 Slides\\
Wellen-Report & Nach jeder Welle & BFE-PL, Fachbereiche & PDF aus Cockpit\\
\rowcolor{falightgray}
Sprint-Review & Alle 2 Wochen & Projektteam & Live-Demo\\
\bottomrule
\end{tabular}

\newpage

% =============================================================================
\section{Qualitätssicherung}
% =============================================================================

\subsection{Review Gates}

\begin{tabular}{c l l p{5cm}}
\toprule
\textbf{Gate} & \textbf{Meilenstein} & \textbf{Prüfer} & \textbf{Kriterien}\\
\midrule
\rowcolor{falightgray}
RG4 & MS7 (Roadmap) & BFE-PL &
    \begin{minipage}[t]{5cm}
    \vspace{0.1cm}
    \begin{itemize}[leftmargin=*, nosep, font=\small]
        \item Baseline analysiert
        \item Zielwerte plausibel
        \item Roadmap abgestimmt
    \end{itemize}
    \vspace{0.1cm}
    \end{minipage}\\

RG5 & MS9 (Module) & BFE-FB, Tech &
    \begin{minipage}[t]{5cm}
    \vspace{0.1cm}
    \begin{itemize}[leftmargin=*, nosep, font=\small]
        \item KPI-Logik validiert
        \item Dashboard getestet
        \item Datenqualität OK
    \end{itemize}
    \vspace{0.1cm}
    \end{minipage}\\

\rowcolor{falightgray}
RG6 & MS10 (Auto) & Tech, BFE-IT &
    \begin{minipage}[t]{5cm}
    \vspace{0.1cm}
    \begin{itemize}[leftmargin=*, nosep, font=\small]
        \item Performance OK
        \item Fehlerbehandlung robust
        \item Doku vollständig
    \end{itemize}
    \vspace{0.1cm}
    \end{minipage}\\

RG7 & MS12 (Abschluss) & \textbf{BFE-GL} &
    \begin{minipage}[t]{5cm}
    \vspace{0.1cm}
    \begin{itemize}[leftmargin=*, nosep, font=\small]
        \item Erfolgskriterien geprüft
        \item Wissenstransfer OK
        \item Regelbetrieb bereit
    \end{itemize}
    \vspace{0.1cm}
    \end{minipage}\\

\bottomrule
\end{tabular}

\newpage

% =============================================================================
\section{EBF-Integration (erweitert)}
% =============================================================================

\subsection{Referenz Phase 1}

\begin{tcolorbox}[infobox={EBF-Grundlagen aus Phase 1}]
\begin{itemize}[leftmargin=*]
    \item \textbf{Messmodell:} A-W-I-T Framework mit EBF-Parametern
    \item \textbf{BCJ-Profil:} Baseline-Daten der Zielgruppe
    \item \textbf{Validierte Parameter:} $\beta$, $\lambda$, $\theta$, $W_{base}$ aus Erstmessung
\end{itemize}
\end{tcolorbox}

\subsection{Erweiterungen Phase 2}

\subsubsection{Interventions-10C-Mapping}

Jede Intervention wird einer 10C-Dimension zugeordnet:

\begin{tabular}{l l l}
\toprule
\textbf{Intervention} & \textbf{10C-Dimension} & \textbf{Ziel}\\
\midrule
\rowcolor{falightgray}
Heizungs-Rechner & AWARE & $A(\text{Heizungsersatz}) \uparrow 0.15$\\
Förder-Konfigurator & WHAT (F) & $u_F(\text{Subvention}) \uparrow$\\
\rowcolor{falightgray}
Nachbarschafts-Vergleich & WHAT (S) & $u_S(\text{Norm}) \uparrow 0.20$\\
\bottomrule
\end{tabular}

\subsubsection{Kampagnen-BCJ-Profile}

Jede neue Kampagne erhält ein eigenes BCJ-Profil:

\begin{tabular}{l l}
\toprule
\textbf{Kampagne} & \textbf{BCJ-Fokus}\\
\midrule
\rowcolor{falightgray}
Heizungsersatz & CONSIDERING $\rightarrow$ INTENDING\\
Solarenergie & INTENDING $\rightarrow$ ACTING\\
\bottomrule
\end{tabular}

\subsubsection{Longitudinale Parameter-Schätzung}

Parameter-Updates basierend auf 4 Wellen:

\begin{tabular}{l c c c}
\toprule
\textbf{Parameter} & \textbf{Welle 1} & \textbf{Ziel Welle 4} & \textbf{Intervention}\\
\midrule
\rowcolor{falightgray}
$\beta$ (present bias) & 0.82 & 0.88 & Commitment-Device\\
$\lambda$ (loss aversion) & 2.1 & -- & Loss-Frame Messaging\\
\rowcolor{falightgray}
$\theta$ (activation) & 0.45 & 0.35 & Reminder-System\\
\bottomrule
\end{tabular}

\newpage

% =============================================================================
\section{Ausblick Phase 3}
% =============================================================================

\begin{tcolorbox}[infobox={Mögliche Ziele Phase 3 (2027)}]
\begin{itemize}[leftmargin=*]
    \item Vollständige Automatisierung aller BFE-Kampagnen
    \item Predictive Analytics (KPI-Prognosen)
    \item Integration mit BFE-CRM und externen Datenquellen
    \item Skalierung auf weitere Bundesämter
\end{itemize}
\end{tcolorbox}

\subsection{Voraussetzungen}

\begin{itemize}[leftmargin=*]
    \item Erfolgreiches Phase-2-Ergebnis ($\geq$80\% Zielerreichung)
    \item Budget-Freigabe durch BFE
    \item Stabiler Regelbetrieb des Cockpits
\end{itemize}

\subsection{Indikative Timeline}

Phase 3: Q1/2027 -- Q4/2027

\vspace{2cm}

\begin{center}
\begin{tcolorbox}[colback=fadarkblue!10, colframe=fadarkblue, width=0.8\textwidth]
\centering
\textbf{\Large Behavioral Impact Cockpit}\\[0.5em]
\textbf{Phase 2: Skalierung \& Optimierung}\\[1em]
\small
Projektplan Version 1.0\\
Erstellt: 02.02.2026\\[1em]
\textit{Aufbauend auf Phase 1 (BFE-2026-001-COCKPIT)}
\end{tcolorbox}
\end{center}

\end{document}
